\section{Pre-processing and Synthetic Generator}
\label{sec:preprocessing}

A \emph{query} is one diagram with one tuple of primitives: segment pairs for \rel{parallel}, \rel{perpendicular} and \rel{equal\_length}, point and circle for \rel{on\_circle}, segment and circle for \rel{tangent}, point and segment for \rel{collinear}. Each diagram contributes at most \queryCap{} queries: \nTrainQueries{} training and \nQueriesId{} / \nQueriesHo{} / \nQueriesEx{} test queries, each a six-way multi-label target with a mask for relations the query kind cannot have. To make the query visible we paint its primitives onto the diagram in green, RGB (\overlayRGB), after resizing the long side to \inputLongSide{} px, a colour absent from the real diagrams. Training images are augmented only in ways that preserve the labels (rotation, flips, scaling, line thickening, noise, JPEG re-encoding); shear would change which lines are parallel.

The synthetic generator exists because the ruler can label anything with coordinates. It samples \synthPointsMin--\synthPointsMax{} points, injects a relation exactly at each construction step with probability \synthInjectRate{}, draws the injected ones as textbook marks, then rotates and jitters every coordinate by \synthJitterPx{} px so that the verifier, not the construction, decides the labels, as for real data. Injected relations pass their check in at least \synthInjectedPassMinPct\% of cases, random candidates in at most \synthRandomPassMaxPct\% (\synthRandomPassEqualLengthPct\% for \rel{equal\_length}). The extrapolation set has \synthExtrapSize{} diagrams of \synthExtrapPointsMin--\synthExtrapPointsMax{} points, larger than anything seen in training.
